\subsection{Exact character correction for every original sieve fibre}
\label{sec:ns-character-sieve}

The torus operator in the preceding subsection acts on the arithmetic
characters of each original shell. We now calculate this action and all
of its fibres, including those needed when its degree shares factors with
the finite character group. No Navier--Stokes existence assertion is a
hypothesis of these calculations.

Fix the original prime $p=12h+1$, $h\geq1$, and an original shell
$3h+1\leq a\leq9h$. Retain $R=4a-p$, $S=pa$, the distinct increasing
prime factors $\ell_1,\ldots,\ell_s$ of $a$, and their original positive
valuations $v_1,\ldots,v_s$. Put
\[
 G=(\mathbb Z/R\mathbb Z)^\times,\qquad
 \widehat G=\operatorname{Hom}(G,\mathbb R/\mathbb Z),\qquad
 \langle\chi,g\rangle=\exp(2\pi i\chi(g)).
\]
We write $G$ multiplicatively and $\widehat G$ additively. Thus
$\chi(gg')=\chi(g)+\chi(g')$ in $\mathbb R/\mathbb Z$, with the
displayed $2\pi i$ retained in every complex character. Let $n=|G|$.
The exact change from the preceding subsection's multiplicative
dual is the isomorphism
$\chi\mapsto(g\mapsto\exp(2\pi i\chi(g)))$; its inverse takes the
unique class in $\mathbb R/\mathbb Z$ with that exponential at each
$g$. The exponential has kernel $\mathbb Z$ on $\mathbb R$ and
maps onto the complex unit circle. Consequently both inverse
compositions are identities and the character equations are preserved.
Thus the two dual conventions are explicitly related, with singleton
fibres; neither is silently substituted for the other.
Every element $g\in G$ has order dividing $n$: its cyclic subgroup
partitions $G$ into cosets, each with that order. Consequently
$\chi(g)\in n^{-1}\mathbb Z/\mathbb Z$. The complete set $\widehat G$
is therefore finite and can be obtained by enumerating all functions
$G\to n^{-1}\mathbb Z/\mathbb Z$ and retaining exactly the identity
and product equations. This is a finite construction on the actual
unit multiplication table; it imposes no cyclicity hypothesis on $G$.

For completeness, these characters separate $G$. If $g$ has order
$o>1$, define $g^j\mapsto j/o$ on its cyclic subgroup. To extend an
additive character from a subgroup $H$ of any finite abelian group to
$H+\mathbb Zx$, choose the least $m>0$ with $mx\in H$, choose any
$\beta\in\mathbb R/\mathbb Z$ with $m\beta=\chi(mx)$, and set
$\chi'(h+jx)=\chi(h)+j\beta$. If $h+jx=h'+j'x$, then
$j-j'=qm$ and $h-h'=-qmx$, so the two values agree. The character
equations follow by addition. Such a $\beta$ exists: if $r\in\mathbb R$
represents $\chi(mx)$, then $r/m$ is a choice. Adjoining elements
until the finite group is exhausted extends the prescribed nonzero
value $1/o$. This proves separation without an unproved extension
assumption.

\begin{lemma}[The original labelled packet as an exact character projection]
\label{lem:ns-labelled-character-packet}
Let $\mathbf X=(X_1,\ldots,X_s)$ be independent formal variables. For
$t\in G$ define
\[
 \begin{split}
 P_a(\chi;\mathbf X)
   &=\prod_{j=1}^s\left(\sum_{e_j=0}^{2v_j}
           X_j^{e_j}\langle\chi,[\ell_j]\rangle^{e_j}\right),\\
 C_{a,t}(\mathbf X)
   &=\frac1{|\widehat G|}\sum_{\chi\in\widehat G}
          \langle\chi,t\rangle^{-1}P_a(\chi;\mathbf X).
 \end{split}
\]
Then the complete coefficient identity is
\[
 C_{a,t}(\mathbf X)=
 \sum_{\substack{0\leq e_j\leq2v_j\ (1\leq j\leq s)\\
                   \prod_j\ell_j^{e_j}\equiv t\pmod R}}
       \prod_{j=1}^s X_j^{e_j}.
\]
In particular every retained coefficient is $1$, and the exponent of
$X_j$ is the original valuation of the corresponding positive divisor
$u=\prod_j\ell_j^{e_j}$. The cardinalities of the original exterior
and middle fibres are respectively
\[
 C_{a,-4^{-1}}(\mathbf1)=|E_a|,\qquad
 C_{a,-a}(\mathbf1)=|M_a|.
\]
\end{lemma}
\begin{proof}
For $g=1$, every character value is $1$. For $g\ne1$, choose
$\eta$ with $\langle\eta,g\rangle\ne1$ by separation. Translation
$\chi\mapsto\chi+\eta$ permutes the finite set $\widehat G$; hence
its character sum equals itself multiplied by
$\langle\eta,g\rangle$ and is zero. Expand the finite product in
$P_a$ without changing any exponent interval. For the term with
the specified original vector $e$, character averaging is therefore
$1$ exactly when $(\prod_j[\ell_j]^{e_j})t^{-1}=1$, and is otherwise
$0$. Distinct original vectors have distinct formal monomials and
distinct positive integers by unique prime factorization. This proves
the identity and both inverse coordinate assertions. Since $R$ is
odd and $\gcd(a,R)=1$, both targets are units. Their respective
equations are exactly $R\mid4u+1$ and $R\mid u+a$.
\end{proof}

The exact comparison with the source's two-dimensional auxiliary torus
is, for each retained unit $g\in G$, the homomorphism
\[
 E_g:\widehat G^2\longrightarrow(\mathbb R/\mathbb Z)^2,
 \qquad E_g(\alpha,\beta)=(\alpha(g),\beta(g)).
\]
The family $(E_g)_{g\in G}$ is injective: equality on all $g$ is
equality of the two functions. Its image consists precisely of pairs
of tables satisfying the character identity and product equations
above. Its inverse reads those two tables. Thus all these coordinates,
not a selected single evaluation, retain the full arithmetic datum.
With the literal source matrix $J_g=\left(\begin{smallmatrix}3&1\\1&5\end{smallmatrix}\right)$,
define
\[
 \mathcal J(\alpha,\beta)=(3\alpha+\beta,\alpha+5\beta).
\]
For every $g$ we have $E_g\mathcal J=J_gE_g$ in the torus, by
evaluating the two displayed sums. This proves the coordinate
morphism to the source operator on every original unit, including
each original prime residue and both targets.

\begin{theorem}[Complete finite image, fibres and correction cosets]
\label{thm:ns-character-coset-correction}
Write $14\widehat G=\{14\chi:\chi\in\widehat G\}$ and
$\widehat G[14]=\{\chi:14\chi=0\}$, and set
$\delta=|\widehat G[14]|$. Then
\[
 \begin{split}
 \ker\mathcal J&=\{(\tau,-3\tau):\tau\in\widehat G[14]\},\\
 \operatorname{im}\mathcal J
   &=\{(\gamma,\eta):\eta-5\gamma\in14\widehat G\}.
 \end{split}
\]
There is an exact sequence of finite abelian groups
\[
 0\longrightarrow\ker\mathcal J\longrightarrow\widehat G^2
   \xrightarrow{\mathcal J}\widehat G^2
   \xrightarrow{q}\widehat G/14\widehat G\longrightarrow0,
 \qquad q(\gamma,\eta)=[\eta-5\gamma].
\]
For every image point, its full inverse fibre is
\[
 \{(\alpha,\gamma-3\alpha):14\alpha=5\gamma-\eta\}.
\]
The displayed root equation is solved by enumerating the already
defined finite set $\widehat G$; all of its solutions are one
solution plus $\widehat G[14]$.

Choose a complete set $\mathcal R\subset\widehat G$ of coset
representatives modulo $14\widehat G$, retaining the representative
table. There are exactly $\delta$ representatives. Define
\[
 T:\widehat G^2\times\mathcal R\longrightarrow\widehat G^2,
 \qquad T(\alpha,\beta,\rho)
       =(3\alpha+\beta,\alpha+5\beta+\rho).
\]
This map is surjective, and every fibre has exactly $\delta$
points. For an arbitrary complex-valued function $F$ on
$\widehat G^2$, or a function with values in a complex polynomial
algebra, the exact identity is
\[
 \frac1{\delta|\widehat G|^2}
 \sum_{\rho\in\mathcal R}\sum_{\alpha,\beta\in\widehat G}
       F(T(\alpha,\beta,\rho))
 =\frac1{|\widehat G|^2}\sum_{\gamma,\eta\in\widehat G}F(\gamma,\eta).
\]
\end{theorem}
\begin{proof}
The first output equation forces $\beta=\gamma-3\alpha$.
Substituting in the second gives $14\alpha=5\gamma-\eta$.
This proves the image, the full fibre and the kernel, including
both directions. The quotient map $q$ is surjective because
$q(0,\eta)=[\eta]$; its kernel is precisely that image. Each fibre
of multiplication by $14$ on $\widehat G$ is a coset of
$\widehat G[14]$, so $|14\widehat G|=|\widehat G|/\delta$ and
$|\widehat G/14\widehat G|=\delta$.

For an arbitrary output $(\gamma,\eta)$ of $T$, choose the unique
recorded representative $\rho$ of $[\eta-5\gamma]$. Then
$14\alpha=5\gamma-\eta+\rho$ has exactly $\delta$ solutions,
and for each of them $\beta=\gamma-3\alpha$. This supplies every
preimage and proves both uniqueness of the coset coordinate and
the full fibre formula. Counting each output once for each of its
$\delta$ preimages proves the averaging identity, coefficient by
coefficient for polynomial-valued $F$.
\end{proof}

\begin{theorem}[Exact alias terms and their removal on original valuations]
\label{thm:ns-sieve-alias-removal}
For $t_1,t_2\in G$ and independent variable tuples $\mathbf X,
\mathbf Z$, put
\[
 F(\gamma,\eta)=
  \langle\gamma,t_1\rangle^{-1}P_a(\gamma;\mathbf X)
  \langle\eta,t_2\rangle^{-1}P_a(\eta;\mathbf Z),
 \qquad G[14]=\{g\in G:g^{14}=1\}.
\]
The uncorrected source matrix gives exactly
\[
 \frac1{|\widehat G|^2}\sum_{\alpha,\beta}
          F(\mathcal J(\alpha,\beta))
   =\sum_{g\in G[14]}
       C_{a,t_1g}(\mathbf X)C_{a,t_2g^{-3}}(\mathbf Z).
\]
The complete correction is
\[
 \frac1{\delta|\widehat G|^2}
     \sum_{\rho\in\mathcal R}\sum_{\alpha,\beta}
          F(T(\alpha,\beta,\rho))
       = C_{a,t_1}(\mathbf X)C_{a,t_2}(\mathbf Z).
\]
Every monomial on the right retains both original exponent vectors.
The corrected identity therefore preserves their complete joint
fibre, not only its cardinality.
\end{theorem}
\begin{proof}
For fixed original exponent vectors $e,f$, write
$u=\prod_j[\ell_j]^{e_j}$, $v=\prod_j[\ell_j]^{f_j}$,
$g=ut_1^{-1}$ and $k=vt_2^{-1}$. Its complex factor under the
uncorrected map is
\[
 \langle\alpha,g^3k\rangle\langle\beta,gk^5\rangle.
\]
The two independent character sums retain the monomial exactly when
$g^3k=1$ and $gk^5=1$. The first equality gives $k=g^{-3}$;
the second then gives $g^{-14}=1$. Conversely those conditions
give both equalities. This is precisely the first displayed sum,
including every original exponent bound and variable.

In the corrected map the same term has the additional factor
$\langle\rho,k\rangle$. For $k^{14}=1$, this value depends only
on the coset of $\rho$ modulo $14\widehat G$, since
$\langle14\chi,k\rangle=\langle\chi,k^{14}\rangle=1$.
If $k\ne1$, separation supplies a character whose value at $k$ is not $1$;
translation by its coset permutes all quotient cosets and forces
the sum over $\mathcal R$ to be zero. If $k=1$, every value is
$1$. In the retained uncorrected terms, $k=g^{-3}$ and $g^{14}=1$.
The two equations $g^3=1$ and $g^{14}=1$ imply $g=1$, since
$1=5\cdot3-14$. Thus precisely $g=k=1$ remains. These are the
original two target equations, proving the second identity with
their full monomials. It also follows directly by applying the
preceding theorem and then the independent character projections.
\end{proof}

\begin{corollary}[All finite covering histories preserve the full ES test]
\label{cor:ns-finite-history-sieve}
Fix an integer $k\geq0$, the original group, and the representative
table $\mathcal R$. For $z_0\in\widehat G^2$ and a retained
sequence $(\rho_0,\ldots,\rho_{k-1})\in\mathcal R^k$, define
\[
 z_{j+1}=\mathcal J z_j+(0,\rho_j)\quad(0\leq j<k),\qquad
 z_k=\mathcal J^k z_0+
       \sum_{j=0}^{k-1}\mathcal J^{k-1-j}(0,\rho_j).
\]
The history-to-output map has precisely $\delta^k$ preimages at
each output, and
\[
 \frac1{\delta^k|\widehat G|^2}
   \sum_{z_0\in\widehat G^2}\sum_{\rho_0,\ldots,\rho_{k-1}\in\mathcal R}
      F(z_k)
       =\frac1{|\widehat G|^2}\sum_zF(z)
\]
for every polynomial-valued $F$. In particular, taking
$F(\gamma,\eta)=\langle\gamma,t\rangle^{-1}P_a(\gamma;\mathbf X)$
recovers exactly $C_{a,t}(\mathbf X)$ for both original targets
after every finite history length.
\end{corollary}
\begin{proof}
The stated expansion of $z_k$ follows by substitution, with the
empty sum and identity map for $k=0$. For a fixed terminal state,
the preceding theorem supplies its unique last coset coordinate
and every one of its $\delta$ predecessor states. Repeating the
same inverse construction at each predecessor gives all histories
and precisely $\delta^k$ of them, with no choices identified.
The averaging formula follows by counting this exact fibre.
For the specified function, the average over $\eta$ contributes
exactly the factor $|\widehat G|$, and the other average is the
proved labelled coefficient projection.
\end{proof}

Returning to integers uses each retained monomial's unique divisor
$u=\prod_j\ell_j^{e_j}$, and then the complete ordered inverse of
Theorem~\ref{thm:shellwise-no-hit}. Its exterior orientation coordinate
and both middle divisors remain part of that return. In particular,
the corrected averages give the same $|E_a|$, $|M_a|$, the same
$q_{a,p}=|E_a|+|M_a|/2$, and the same $Q_p$ and $\mathcal O_p$
in Theorem~\ref{thm:global-no-hit}. This supplies an exact arithmetic
consequence of the source cover, including a proved obstruction to
using its uncorrected finite sampling.

The source's smooth fields use continuous auxiliary coordinates.
Here $E_g$ is the explicit finite-coordinate map to that torus, and
the preceding exact sequences and character products describe its
arithmetic restriction. The operator correspondence by itself does
not construct a Navier--Stokes velocity or force from a prime, or
establish nonemptiness of any previously empty ES fibre. No such
construction or prime-coverage conclusion is a premise of a result
above.
